% =============================================================================
% UQFF v5.78  Section Template :: T-LAG
% Lagrangian-derivation papers (action-principle / Euler-Lagrange / EoM)
% -----------------------------------------------------------------------------
% USAGE
%   1. Copy to whitepapers/PAPER_<NNNN>_<slug>.tex
%   2. Replace every <PLACEHOLDER>
%   3. Keep CLOSURE / TEMPLATE / CANONICAL banner lines VERBATIM
%   4. pdflatex direct, two passes.
% -----------------------------------------------------------------------------
% Canonical Lagrangian skeleton (UQFF v5.78):
%   L_UQFF = L_kinetic + L_grav + L_buoy + L_resonance - V(rho_vac)
%   with the buoyancy term L_buoy = (1/2) * rho_vac * (4 pi / 3) * c^2 * phi^2 * cos(pi t_n)
%   yielding the Universal Inertia invariant U_I = -dL_buoy / d phi^2 = ...
% =============================================================================
\documentclass[11pt]{article}
\usepackage[margin=1in]{geometry}
\usepackage{amsmath,amssymb,amsthm}
\usepackage{booktabs,longtable,array}
\usepackage{listings,xcolor}
\usepackage[colorlinks=true,linkcolor=blue,urlcolor=blue,citecolor=blue]{hyperref}
\lstset{basicstyle=\ttfamily\small,breaklines=true,frame=single,
        language=Python,showstringspaces=false,columns=fullflexible}

\title{<PAPER TITLE>: \\ Lagrangian Derivation under UQFF v5.78}
\author{Star-Magic UQFF \\ PAPER\_<NNNN> \,/\, Session <SESSION>}
\date{<DATE>}

% CLOSURE :: <closure_label> :: predicted=<X> observed=<Y> error_pct=<Z>
% TEMPLATE :: T-LAG
% CANONICAL :: lagrangian_skeleton = L_kin + L_grav + L_buoy + L_res - V(rho_vac)

\begin{document}
\maketitle

\begin{abstract}
<State the field / interaction whose Lagrangian is being derived,
the canonical UQFF skeleton that hosts it, the Euler--Lagrange equation(s)
obtained, and the observable they predict.>
\end{abstract}

\section{Field Content}
\begin{itemize}
  \item Scalar field(s): $\phi_{\mathrm{SCm}}$, $\phi_{\mathrm{UA}}$ -- vacuum order parameters.
  \item Vector/tensor field(s): <list, with mass dimension>.
  \item Coupling constants: $\beta_i \approx 0.603$, $\kappa = 5\times 10^{-4}/\mathrm{day}$, $F_{TRZ}=0.1$.
\end{itemize}

\section{Lagrangian Density}
\begin{align}
  \mathcal{L}_{\mathrm{UQFF}} &= \mathcal{L}_{\mathrm{kin}} + \mathcal{L}_{\mathrm{grav}} + \mathcal{L}_{\mathrm{buoy}} + \mathcal{L}_{\mathrm{res}} - V(\rho_{\mathrm{vac}}) \\
  \mathcal{L}_{\mathrm{buoy}} &= \tfrac{1}{2}\,\rho_{\mathrm{vac}}\,\tfrac{4\pi}{3}\,c^{2}\,\phi^{2}\cos(\pi t_n) \label{eq:Lbuoy}
\end{align}
with $\rho_{\mathrm{vac}}\in\{\rho_{\mathrm{vac,SCm}},\rho_{\mathrm{vac,UA}}\}$ depending on the regime
(structural G9 vs G7, see Sec.~\ref{sec:closure}).

\section{Euler--Lagrange Equation}
\begin{equation}
  \partial_\mu\!\left(\frac{\partial\mathcal{L}}{\partial(\partial_\mu\phi)}\right) - \frac{\partial\mathcal{L}}{\partial\phi} = 0
  \;\;\Longrightarrow\;\;
  \boxed{\;<EoM equation here>\;}
  \label{eq:tlag-eom}
\end{equation}

\section{Universal Inertia Invariant}
The invariant differential arising from \eqref{eq:Lbuoy}:
\begin{equation}
  U_I \;=\; 3\,\rho_{\mathrm{vac}}\,\tfrac{4\pi}{3}\,c^{2}\,\cos(\pi t_n)
\end{equation}
Sign regimes: $U_I>0$ centripetal (mass-building), $U_I=0$ zero-point gravity
($t_n = \tfrac{1}{2}, \tfrac{3}{2}, \ldots$), $U_I<0$ centrifugal (negentropic void).

\section{Observable Prediction \& Closure}
\label{sec:closure}
\begin{longtable}{lll}
\toprule
Quantity & Predicted (UQFF v5.78) & Observed (cite) \\
\midrule
<observable> & <value> & <value> [\cite{<ref>}] \\
\bottomrule
\end{longtable}

\section{Cross-references}
\begin{itemize}
  \item Closed-form $r_{hz}$: \texttt{QCalcGeom.py::solve\_habitable\_zone} (v2.3.0).
  \item Universal Inertia: \texttt{QCalcGeom.py::universal\_inertia}.
  \item Parent skeleton: \texttt{uqff\_lagrangian\_derivation.py}.
\end{itemize}

\bibliographystyle{plain}
\begin{thebibliography}{9}
\bibitem{ref1} <DOI/arXiv real citation>.
\bibitem{paper409}  PAPER\_409 (Star-Magic): hydrogen geometry $\rho_{\mathrm{vac,SCm}}$.
\bibitem{paper1160} PAPER\_1160 (Star-Magic): G7 $|SO(5)|$ closure.
\end{thebibliography}

\end{document}
